\begin{table}[htbp]
    \centering
    \small\setlength{\tabcolsep}{5pt}
    \begin{tabular}{lrrrrr}
    \hline
    Dataset & $p$ & Effective rank & $p/\mathrm{eff}$ & Mean VIF & $\Delta$ AUC (LASSO) \\
    \hline
    \texttt{ind5} & 51 & 11.9 & 4.3 & 8.9 & $+0.0249$$^{*}$ \\
    \texttt{ind10} & 101 & 19.3 & 5.2 & 10.8 & $+0.0396$$^{*}$ \\
    \texttt{ind12} & 121 & 19.5 & 6.2 & 13.6 & $+0.0306$$^{*}$ \\
    \texttt{ind17} & 171 & 23.5 & 7.3 & 13.8 & $+0.0602$$^{*}$ \\
    \texttt{ind30} & 301 & 32.3 & 9.3 & 18.5 & $+0.1210$$^{*}$ \\
    \texttt{ind38} & 361 & 36.6 & 9.9 & 20.4 & $+0.0948$$^{*}$ \\
    \texttt{ind48} & 481 & 42.8 & 11.2 & 28.2 & $+0.0604$$^{*}$ \\
    \texttt{ff49ind} & 491 & 42.9 & 11.4 & 29.0 & $+0.0822$$^{*}$ \\
    \texttt{sector\_etf} & 58 & 13.3 & 4.4 & 14.9 & -- \\
    \texttt{fama\_french} & 50 & 20.9 & 2.4 & 3.6 & -- \\
    \hline
    \end{tabular}
    \caption{Collinearity across the ladder. Effective rank is $\exp(H)$ for the Shannon entropy $H$ of the normalised eigenvalue spectrum of the feature correlation matrix; $p/\mathrm{eff}$ is the resulting compression factor. Mean VIF is the mean of $\mathrm{diag}(R^{-1})$, computed by pseudo-inverse because these matrices are severely ill-conditioned at the wider widths --- where the two measures disagree, the effective rank is the more stable and should be preferred. The final column is the logistic-regression gain for LASSO at that width, from Table~\ref{tab:dimreal}. \texttt{sector\_etf} and \texttt{fama\_french} are included for context and are not part of the ladder.}
    \label{tab:collinearity}
\end{table}
